\documentclass[11pt, a4paper]{article}

\usepackage[top=2cm, bottom=2cm, left=2.5cm, right=2cm]{geometry}
\usepackage{fontspec}
\usepackage{amsmath, amssymb}
\usepackage{booktabs}
\usepackage{tabularx}
\usepackage{graphicx}
\usepackage{float}
\usepackage{xcolor}
\usepackage{hyperref}
\usepackage{xurl}
\usepackage{parskip}
\usepackage{listings}
\usepackage{microtype}

\graphicspath{{images/}}

\hypersetup{
    colorlinks=true,
    linkcolor=blue!70!black,
    citecolor=blue!70!black,
    urlcolor=blue!70!black
}

\lstset{
    basicstyle=\small\ttfamily,
    breaklines=true,
    frame=single,
    backgroundcolor=\color{gray!10}
}

\begin{document}

% ---------------------------------------------------------
% COVER PAGE
% ---------------------------------------------------------
\begin{titlepage}
\begin{center}

\vspace*{0.5cm}

\Large
\textbf{POSTS AND TELECOMMUNICATIONS INSTITUTE OF TECHNOLOGY}\\
\textbf{Faculty of Information Technology}

\vspace{2.0cm}

\Huge
\textbf{Intelligent System Development}\\
\vspace{0.4cm}
\LARGE
\textbf{ASSIGNMENT 03}\\
\vspace{0.6cm}
\Large
\textbf{From Data Representation to Deep Learning}

\vspace{1.2cm}

\large
\textit{"Deep Learning = Function Composition + Representation Learning + Optimization."}

\vfill

\normalsize
\begin{tabular}{rl}
\textbf{Instructor / Lecturer:} & Assoc. Prof. Dinh Que Tran, Ph.D. \\
\textbf{Email:} & quetd@ptit.edu.vn, tdque@yahoo.com \\[0.8cm]
\textbf{Student Name:} & Lưu Anh Dũng \\
\textbf{Student ID:} & B23DCDK036 \\
\textbf{Semester:} & I.2025 -- 2026 \\
\textbf{Class:} & E23CNPM02 \\
\textbf{Submission:} & Individual Project \\
\end{tabular}

\vfill

\textbf{Hanoi, 2026}

\vspace{0.5cm}

\end{center}
\end{titlepage}

% ---------------------------------------------------------
% TABLE OF CONTENTS
% ---------------------------------------------------------
\tableofcontents
\newpage

% ---------------------------------------------------------
% SECTION 1: EXECUTIVE SUMMARY
% ---------------------------------------------------------
\section{Executive Summary}

This report extends the three intelligent systems from Assignment 02 -- diabetes risk
screening, house-price regression, and e-commerce review classification -- with a deep
learning model implemented entirely from scratch using NumPy, with no TensorFlow, PyTorch,
or Keras anywhere in the implementation. For each application, four models are trained on
an identical train/test split and compared: three classical scikit-learn models and one
from-scratch feedforward network whose forward propagation, loss, backpropagation, and
gradient-descent update are all derived and coded by hand via the chain rule.

Each dataset was also grown relative to Assignment 02, per the separate verbal requirement
that every assignment's dataset be larger than the one before it, while staying well under
a gigabyte in total.

\begin{table}[H]
\centering
\begin{tabular}{lll}
\toprule
\textbf{Application} & \textbf{Task} & \textbf{Dataset (rows)} \\
\midrule
Diabetes & Binary classification & BRFSS 2019+2020 (815,711 raw) \\
House Price & Regression & USA real estate, 600K sample (2,226,382 raw) \\
Customer Behaviour & Binary classification & Flipkart reviews (363,261 raw) \\
\bottomrule
\end{tabular}
\caption{Overview of the Three Intelligent Systems}
\end{table}

\textbf{1. Diabetes Prediction}
\begin{itemize}
    \item \textbf{Dataset:} CDC BRFSS 2019 + 2020 (concatenated), 815,711 raw rows.
    \item \textbf{Problem:} Binary classification -- predict whether a survey respondent is
    diabetic or pre-diabetic.
    \item \textbf{Representation:} Scaled numeric + one-hot nominal + binary-passthrough
    feature matrix, $X \in \mathbb{R}^{815,711 \times 52}$.
    \item \textbf{4 models:} Logistic Regression, Random Forest, Gradient Boosting,
    from-scratch DL ($52 \to 32 \to 16 \to 1$).
    \item \textbf{Winner (Accuracy/F1):} Gradient Boosting (Accuracy 0.852, F1 0.286); the
    from-scratch DL model trails badly on Recall (0.028) due to unweighted BCE on an
    imbalanced target.
\end{itemize}

\textbf{2. House Price Prediction}
\begin{itemize}
    \item \textbf{Dataset:} USA Real Estate Listings, 600,000-row sample of 2,226,382 raw
    listings.
    \item \textbf{Problem:} Regression -- predict listing price ($\log(1+y)$ in USD).
    \item \textbf{Representation:} Scaled numeric + one-hot categorical + frequency-encoded
    city feature matrix, $X \in \mathbb{R}^{599,522 \times 65}$.
    \item \textbf{4 models:} Linear Regression, Random Forest, Gradient Boosting,
    from-scratch DL ($65 \to 64 \to 32 \to 1$, linear output).
    \item \textbf{Winner ($R^2$):} Random Forest ($R^2$ 0.417, RMSE \$794,082); the
    from-scratch DL model places last ($R^2$ 0.057).
\end{itemize}

\textbf{3. E-Commerce Customer Behaviour}
\begin{itemize}
    \item \textbf{Dataset:} Flipkart product reviews, 363,261 raw rows (different platform
    from Assignment 02's Sephora reviews).
    \item \textbf{Problem:} Binary classification -- predict whether a review is a
    "recommend" (\texttt{Rate} $\geq 4$).
    \item \textbf{Representation:} TF-IDF (1--2 gram) over review text + scaled price;
    classical models use the full 8,695-term vocabulary, the DL model a 2,000-term cap.
    \item \textbf{4 models:} Logistic Regression, Linear SVM, Multinomial Naive Bayes,
    from-scratch DL ($2{,}001 \to 128 \to 32 \to 1$).
    \item \textbf{Winner (Accuracy/F1):} Logistic Regression (Accuracy 0.994, F1 0.996);
    the from-scratch DL model collapses to a majority-class predictor (Accuracy 0.766,
    Recall 1.0, Precision = Accuracy).
\end{itemize}

Two findings stand out across all three applications. First, the from-scratch deep
learning model never wins its task's primary metric -- it places last (or effectively
last) in every one of the three applications. Second, the same underlying cause explains
two of the three shortfalls: the from-scratch training loop uses plain, unweighted
full-batch gradient descent with no class weighting, no resampling, and a fixed 0.5
threshold, which is vulnerable to exactly the class imbalance present in both
classification tasks (84.4\%/15.6\% for diabetes, 76.5\%/23.5\% for customer behaviour) in
a way the classical models' training procedures are not. The house-price shortfall has a
different, complementary cause: a heavy-tailed regression target colliding with an
unbounded linear output layer that lacks an ensemble's implicit averaging/regularization.
All three are genuine, explainable limitations of the from-scratch implementation as
built, not implementation bugs -- every backpropagation derivation was separately
gradient-checked against numerical finite differences before being used to train on real
data (see each application's DL subsection).

% ---------------------------------------------------------
% SECTION 2: CONNECTION TO LECTURE 03
% ---------------------------------------------------------
\section{Connection to Lecture 03: Representation Learning and Deep Learning}
\label{sec:lecture}

This section is written once, at the framework level, so the three per-application DL
subsections below do not each have to re-derive the same mathematics -- each one
instantiates the ideas here with its own concrete widths and numbers, and references this
section rather than repeating it.

\subsection{The central equation}

Lecture 03 frames a deep neural network as a composition of simple functions applied one
after another to the input:
\[
\hat{y} = (f_L \circ f_{L-1} \circ \cdots \circ f_1)(x)
\]
Each $f_\ell$ is one layer: a linear transformation followed by a nonlinear activation.
Stacking $L$ of these layers is what makes the network "deep." The lecture's own closing
summary states this compactly: \textit{Deep Learning = Function Composition +
Representation Learning + Optimization}. Function composition is the equation above;
representation learning is the idea that each layer's output is a new, learned
representation of the input, progressively more useful for the final task than the raw
features were; optimization is the process -- gradient descent -- that adjusts every
layer's weights so the composed function actually produces useful representations rather
than arbitrary ones.

\subsection{Forward propagation}

A forward pass pushes a batch of input rows through every layer in sequence, always in the
same two-step order: multiply first, activate second.
\[
Z_\ell = A_{\ell-1} W_\ell + b_\ell, \qquad A_\ell = g_\ell(Z_\ell)
\]
$A_0 = X$ is the input itself, $W_\ell$/$b_\ell$ are that layer's learned weights and
biases, $Z_\ell$ is the pre-activation, and $g_\ell$ is the layer's activation: ReLU for
every hidden layer across all three applications, and either sigmoid (binary
classification) or no activation at all (regression) at the output layer. Every
intermediate value computed during the forward pass is cached, because the backward pass
immediately after it needs those exact values to compute gradients.

\subsection{Loss}

The loss is the single number that measures how wrong the network's current predictions
are. Binary classification uses binary cross-entropy,
$-\text{mean}(y\log\hat{y} + (1-y)\log(1-\hat{y}))$, which reduces to
$-\log(\text{the probability assigned to the correct class})$ per row; regression uses mean
squared error, $\text{mean}((\hat{y}-y)^2)$, which penalizes large misses
disproportionately more than small ones because the error is squared before averaging.

\subsection{Backpropagation}

Backpropagation answers one question per layer, working backward from the output: how much
is this layer to blame for the final error, and in which direction should its weights move
to reduce it? This is the chain rule applied by hand, one layer at a time. For the output
layer, the error signal is $dZ_L = \hat{y} - y$ (classification, sigmoid+BCE) or
$dZ_L = 2(\hat{y}-y)/n$ (regression, linear+MSE) -- different formulas, because
classification's simple form is a known algebraic cancellation between sigmoid's derivative
and BCE's derivative that does not occur for a linear output under MSE. Every earlier layer
then follows the same repeating pattern: route the error backward through the next layer's
weights ($dH_{\ell} = dZ_{\ell+1} W_{\ell+1}^{\top}$), gate it by the ReLU derivative
($dZ_\ell = dH_\ell \odot \mathbb{1}[Z_\ell > 0]$, zeroing the gradient for any neuron that
was inactive during the forward pass), then compute that layer's parameter gradients
($dW_\ell = A_{\ell-1}^{\top} dZ_\ell$, $db_\ell = \sum dZ_\ell$).

\subsection{Gradient descent}

"Learning" is repeating four steps every epoch: \textbf{forward} (compute predictions)
$\to$ \textbf{loss} (measure how wrong they are) $\to$ \textbf{backward} (compute, via the
chain rule above, how much each weight is to blame) $\to$ \textbf{update} (nudge every
weight a small step opposite its gradient, $W \mathrel{-}= \eta \, dW$, since the gradient
points toward higher loss). Every network in this assignment uses full-batch gradient
descent -- one forward/backward pass over the entire training set per epoch, no
mini-batching -- the simplest variant, and the one explicitly described in the lecture
material.

\subsection{Weight initialization}

Every weight matrix is drawn from $\mathcal{N}(0,1)$ and scaled by $\sqrt{2/\text{fan\_in}}$
(He initialization), biases start at zero. Random initialization breaks the symmetry that
would otherwise make every neuron in a layer compute the identical thing; the
$\sqrt{2/\text{fan\_in}}$ scale keeps each layer's output variance roughly stable rather
than exploding or vanishing as it passes through multiple layers of ReLU units.

\subsection{What makes this "from scratch"}

Every one of the six per-layer gradients ($dW_1, db_1, dW_2, db_2, dW_3, db_3$) in every
application's network is derived symbolically via the chain rule and implemented as plain
NumPy array operations -- no autograd engine, no framework. Each application's DL section
below states that its specific backward-pass implementation was separately verified against
a numerical finite-difference gradient check on synthetic data before being used to train on
real data, confirming the hand-derived formulas are correct independent of whether the
resulting model performs well.

% ---------------------------------------------------------
% SECTION 3: APPLICATION 1 - DIABETES
% ---------------------------------------------------------
\section{Application 1 --- Diabetes Prediction}

\subsection{Problem Description}
Binary classification: predict whether a BRFSS survey respondent is diabetic or
pre-diabetic from self-reported lifestyle and biometric indicators, extending Assignment
02's diabetes screening task onto a larger, two-year BRFSS sample.
\[
f: \mathbb{R}^{52} \longrightarrow \{0, 1\}
\]

\subsection{Dataset}
\begin{itemize}
    \item \textbf{Source:} CDC BRFSS, Kaggle dataset
    \texttt{spandanjit2005/brfss-diabetes-indicator-dataset}.
    \item \textbf{Files used:} \texttt{DATASET\_2019.csv} + \texttt{DATASET\_2020.csv},
    concatenated with a \texttt{year} indicator.
    \item \textbf{Rows:} 815,711 raw $\to$ 652,568 train / 163,143 test (80/20 stratified
    split; no rows dropped during cleaning, since all 33 raw columns are retained or
    encoded rather than filtered out).
    \item \textbf{Growth vs.\ Assignment 02:} 3.2$\times$ the prior 253,680-row BRFSS 2015
    sample.
    \item \textbf{Target:} \texttt{diabetes} (raw 3-class text: Non-diabetic / Pre-diabetic
    / Diabetic) collapsed to \texttt{Diabetes\_binary} $\in \{0,1\}$ (0 = confirmed
    non-diabetic, 1 = pre-diabetic or diabetic).
\end{itemize}

\subsection{Data Understanding and Cleaning}
The raw target is 84.4\% Non-diabetic, 13.8\% Diabetic, 2.2\% Pre-diabetic; collapsing the
latter two gives an 84.4\%/15.6\% binary split (Figure~\ref{fig:diabetes-target}). Seven
columns were dropped: five raw-text duplicates of an already-numeric ordinal column (e.g.\
\texttt{general\_health\_00} duplicating \texttt{gen\_health\_00}), and two columns more
than 48\% missing across all 815,711 rows. Two \(\sim\)50\%-missing but clinically
important columns (\texttt{high\_bp\_00}, \texttt{high\_cholesterol\_00}) were kept rather
than dropped, each paired with an explicit \texttt{\_missing} indicator column recording
whether the value was present at all, before median/mode imputation (fit on the training
split only) filled the gap. \texttt{year} was tested rather than assumed useful: the
positive-diagnosis rate is 15.97\% in 2019 vs.\ 15.20\% in 2020
(Figure~\ref{fig:diabetes-year}) -- a modest but real \(\sim\)0.8-point gap -- so it was
kept as a single \texttt{is\_2020} binary flag.

\begin{figure}[H]
\centering
\includegraphics[width=0.68\textwidth]{diabetes_target_balance.png}
\caption{Target class balance, raw 3-class (before collapsing to binary)}
\label{fig:diabetes-target}
\end{figure}

\begin{figure}[H]
\centering
\includegraphics[width=0.68\textwidth]{diabetes_bmi_by_target.png}
\caption{BMI by diabetes status. Mean BMI: 27.7 (non-diabetic) vs.\ 31.6 (diabetic) --
the clinically expected separation.}
\label{fig:diabetes-bmi}
\end{figure}

\begin{figure}[H]
\centering
\includegraphics[width=0.68\textwidth]{diabetes_genhealth_by_target.png}
\caption{Self-rated general health (1=Excellent..5=Poor) by diabetes status. Mean: 2.40
(non-diabetic) vs.\ 3.22 (diabetic) -- diabetics report meaningfully worse general health.}
\label{fig:diabetes-genhealth}
\end{figure}

\begin{figure}[H]
\centering
\includegraphics[width=0.6\textwidth]{diabetes_year_rate.png}
\caption{Positive-diagnosis rate by survey year (2019 vs.\ 2020)}
\label{fig:diabetes-year}
\end{figure}

\begin{figure}[H]
\centering
\includegraphics[width=0.68\textwidth]{diabetes_corr_heatmap.png}
\caption{Correlation heatmap: numeric features + target. \texttt{gen\_health\_00} and
\texttt{bmi\_00} show the strongest correlation with \texttt{Diabetes\_binary} among the
numeric columns available at this stage; no pair of input features exceeds $|r| > 0.7$,
so no column is dropped on multicollinearity grounds.}
\label{fig:diabetes-corr}
\end{figure}

\subsection{Representation}
The final feature matrix combines three blocks, concatenated column-wise after a
train/test split computed \emph{before} any fitting step:
\begin{itemize}
    \item \textbf{Scaled numeric} (\texttt{StandardScaler}, fit on train only): weight,
    height, BMI, physical/mental health days, age, education, income, general health,
    last-checkup recency, plus the two \(\sim\)50\%-missing clinical columns after
    imputation -- 12 columns.
    \item \textbf{Binary passthrough} (mode-imputed on train, unscaled): sex and 7 boolean
    survey fields, plus \texttt{is\_2020} and the two missing-value indicator flags -- 11
    columns.
    \item \textbf{One-hot nominal} (categories fit on train only, unseen test categories
    ignored): race, marital status, employment status, personal-doctor status -- 29
    columns.
\end{itemize}
Resulting shape: $X \in \mathbb{R}^{815,711 \times 52}$.

\subsection{Classical Machine-Learning Models}
\begin{table}[H]
\centering
\begin{tabularx}{\textwidth}{@{}l X@{}}
\toprule
\textbf{Model} & \textbf{Why} \\
\midrule
Logistic Regression & Same functional form as the DL model's output layer alone -- a
useful baseline for what depth adds. \\
Random Forest & Model deployed in Assignment 02 for this app -- kept for continuity. \\
Gradient Boosting & Sequential residual-correction; structurally different from Random
Forest's parallel averaging. \\
\bottomrule
\end{tabularx}
\caption{Diabetes: 3 classical models and selection rationale}
\end{table}

\subsection{From-Scratch Deep Learning Model}
\textbf{Architecture:} $52 \to 32 \to 16 \to 1$, ReLU hidden layers, sigmoid output, binary
cross-entropy loss. $32 \to 16$ follows the lecture tutorial's $8 \to 16 \to 8 \to 1$
pattern scaled up for this wider 52-feature input (first hidden layer $\approx$ input
width, second layer half that).

\textbf{Forward propagation:} layer 1 linearly recombines the 52 scaled/encoded survey
features and applies ReLU, introducing the nonlinearity a purely linear stack could never
have (composing linear functions collapses to one linear function, however many layers are
stacked); layer 2 recombines layer 1's learned features into higher-level ones --
representation learning in miniature, since layer 2 never sees BMI or blood pressure
directly, only whatever combinations of them layer 1 found useful; the output layer applies
one more linear combination then sigmoid, squashing the result to a $(0,1)$ diabetes
probability.

\textbf{Backpropagation:} $dZ_3 = \hat{y}-y$ is the output error signal (this simple form is
sigmoid's and BCE's derivatives canceling algebraically, a known property of this exact
pairing); layer 2 asks how much it contributed to that error by routing the signal back
through $W_3$ ($dH_2 = dZ_3 W_3^{\top}$), then the ReLU derivative zeroes the gradient for
any neuron that was inactive during the forward pass ($dZ_2 = dH_2 \odot
\mathbb{1}[Z_2>0]$); the same route-back$\to$gate$\to$gradient pattern repeats once more for
layer 1. This implementation was separately gradient-checked against a numerical
finite-difference estimate on synthetic data (maximum relative error $\approx 5\times
10^{-8}$) before training on the real BRFSS data.

\textbf{Training:} 200 epochs, full-batch gradient descent, learning rate 0.1. Loss fell
smoothly from 0.679 to 0.353 with no sign of divergence (Figure~\ref{fig:diabetes-loss}).

\begin{figure}[H]
\centering
\includegraphics[width=0.65\textwidth]{diabetes_loss_curve.png}
\caption{From-scratch DL model: training loss (binary cross-entropy) vs.\ epoch}
\label{fig:diabetes-loss}
\end{figure}

\subsection{Evaluation: All 4 Models}
\begin{table}[H]
\centering
\begin{tabular}{lcccc}
\toprule
\textbf{Model} & \textbf{Accuracy} & \textbf{Precision} & \textbf{Recall} & \textbf{F1} \\
\midrule
Logistic Regression & 0.8498 & 0.5563 & 0.1814 & 0.2736 \\
Random Forest & 0.8489 & 0.5494 & 0.1734 & 0.2636 \\
\textbf{Gradient Boosting} & \textbf{0.8516} & \textbf{0.5723} & \textbf{0.1905} & \textbf{0.2859} \\
DL (from scratch) & 0.8450 & 0.5592 & 0.0284 & 0.0541 \\
\bottomrule
\end{tabular}
\caption{Diabetes: 4-model test-set comparison}
\label{tab:diabetes-results}
\end{table}

\begin{figure}[H]
\centering
\includegraphics[width=0.72\textwidth]{diabetes_accuracy_bar.png}
\caption{Accuracy across all 4 models}
\label{fig:diabetes-acc}
\end{figure}

\begin{figure}[H]
\centering
\includegraphics[width=0.8\textwidth]{diabetes_prf1_bar.png}
\caption{Precision / Recall / F1 across all 4 models}
\label{fig:diabetes-prf1}
\end{figure}

\begin{figure}[H]
\centering
\includegraphics[width=0.55\textwidth]{diabetes_confusion_matrix.png}
\caption{From-scratch DL model: confusion matrix on the test set}
\label{fig:diabetes-cm}
\end{figure}

\begin{figure}[H]
\centering
\includegraphics[width=0.72\textwidth]{diabetes_roc_curves.png}
\caption{ROC curves across all 4 models. The from-scratch DL model's AUC (0.8025) sits
close behind the three classical models (0.8098--0.8239), a much smaller gap than the
Accuracy/Recall table above suggests.}
\label{fig:diabetes-roc}
\end{figure}

\begin{figure}[H]
\centering
\includegraphics[width=0.78\textwidth]{diabetes_feature_importance.png}
\caption{Top 15 feature importances, Gradient Boosting (the winning classical model).
\texttt{gen\_health\_00}, \texttt{age\_00}, and \texttt{bmi\_00} -- the same features the
correlation heatmap flagged as most predictive -- dominate.}
\label{fig:diabetes-fi}
\end{figure}

\subsection{Written Comparison}
On Accuracy, Gradient Boosting wins narrowly (0.8516) ahead of Logistic Regression
(0.8498), Random Forest (0.8489), and the from-scratch DL model last (0.8450) -- but
Accuracy alone is misleading on this 84.4\%/15.6\%-imbalanced target, since a trivial
always-negative classifier already scores \(\sim\)84.4\%. Recall is the clinically relevant
metric, and here the gap is stark: Gradient Boosting leads at 0.191, Logistic Regression
0.181, Random Forest 0.173, and the from-scratch DL model trails far behind at just 0.028
(723 true positives caught out of 25,441 actual positives in the test set).

The ROC curves in Figure~\ref{fig:diabetes-roc} resolve what would otherwise look like a
clear failure: the DL model's AUC (0.8025) sits close behind the three classical models
(0.8098--0.8239), not far below them -- meaning the network did learn a genuinely useful
ranking of who is more/less likely diabetic, it is not confused or undertrained. The
Accuracy/Recall gap is specifically a \textbf{decision-threshold} problem, not a
\textbf{ranking} problem: at the fixed 0.5 cutoff, unweighted binary cross-entropy
training on an 84.4\%-negative target pushed the network's probability outputs low
enough that few positives cross 0.5, even though the relative ordering of those
probabilities (what AUC measures) is nearly as good as the tree ensembles'.

In the lecture's representation-capacity/overfitting/dataset-size framing: this is not an
overfitting problem (the loss curve fell smoothly with no divergence) or a
representation-capacity problem ($52 \to 32 \to 16 \to 1$ has ample capacity for 52 tabular
features, and Figure~\ref{fig:diabetes-fi} confirms the features Gradient Boosting relies
on most are exactly the ones the correlation heatmap flagged as predictive, so the
representation itself is sound) -- it is an optimization/class-imbalance problem specific
to plain full-batch gradient descent at a fixed threshold. Gradient Boosting's sequential
residual-correction and Random Forest's per-tree resampling both implicitly re-weight
harder (minority-class) examples over training, while plain gradient descent on the full
batch does not. This is a genuine, explainable limitation of the from-scratch
implementation as built, not a bug -- a class-weighted loss term or a tuned decision
threshold (Figure~\ref{fig:diabetes-roc} shows a lower threshold would recover much of
the missing Recall) are natural extensions, left out here to keep the training loop
matching the lecture's exact unweighted BCE formula.

% ---------------------------------------------------------
% SECTION 4: APPLICATION 2 - HOUSE PRICE
% ---------------------------------------------------------
\section{Application 2 --- House Price Prediction}

\subsection{Problem Description}
Regression: predict a US residential listing's price from structural and geographic
attributes.
\[
f: \mathbb{R}^{65} \longrightarrow \mathbb{R}^{+}
\]
evaluated via RMSE, MAE, and $R^2$ in dollar terms after inverting the $\log(1+y)$
training target.

\subsection{Dataset}
\begin{itemize}
    \item \textbf{Source:} Kaggle dataset
    \texttt{ahmedshahriarsakib/usa-real-estate-dataset}, \texttt{realtor-data.csv}.
    \item \textbf{Rows:} 2,226,382 raw listings, sampled once (seeded,
    \texttt{RANDOM\_SEED=42}) to a persisted 600,000-row working set
    (\texttt{realtor-data-sample.csv}), cleaned to 599,522 rows, split 479,617 train /
    119,905 test.
    \item \textbf{Growth vs.\ Assignment 02:} 2.5$\times$ the prior 238,924-row Vietnamese
    real-estate dataset, and a genuinely different (US) market.
    \item \textbf{Target:} \texttt{price} (USD), trained as $\log(1+\text{price})$ and
    inverted with $\text{expm1}$ for evaluation.
\end{itemize}

\subsection{Data Understanding and Cleaning}
Raw \texttt{price} is heavily right-skewed, with a handful of implausible values including
one listing at \$2,147,483,600 -- a value suspiciously close to $2^{31}-1$, a classic
32-bit integer overflow artifact -- and 280 rows at \$0 or negative
(Figure~\ref{fig:house-price-dist}). Those rows were dropped rather than clipped, since
they are not genuine sale prices. \texttt{house\_size}, \texttt{bed}, \texttt{bath}, and
\texttt{acre\_lot} showed the same overflow-style corruption (\texttt{house\_size} up to
1,040,400,400; \texttt{bed} up to 473; \texttt{bath} up to 830) that drove the raw
\texttt{house\_size}-vs-\texttt{price} correlation to an implausible $r \approx 0.0004$
(Figure~\ref{fig:house-price-scatter}); capping each at its own 99th percentile (rather
than dropping rows) raised the correlation to a sensible $r \approx 0.27$ with raw price
while preserving 99\% of real data. \texttt{brokered\_by} and \texttt{street} were dropped
as anonymized, high-cardinality identifiers with no generalizable signal.
\texttt{prev\_sold\_date} was engineered into \texttt{days\_since\_prev\_sale} (correlation
$\approx 0.036$ with price -- weak but real and free to keep) rather than dropped outright.

\begin{figure}[H]
\centering
\includegraphics[width=0.85\textwidth]{house_price_dist.png}
\caption{Raw price distribution, clipped for display at the 99.5th percentile (left) vs.\
$\log(1+\text{price})$ (right) -- the log-transform normalizes the heavy right tail,
justifying training in log-space}
\label{fig:house-price-dist}
\end{figure}

\begin{figure}[H]
\centering
\includegraphics[width=0.62\textwidth]{house_price_size_scatter.png}
\caption{\texttt{house\_size} vs.\ price, 20K-point sample -- reveals the outlier-driven
distortion before capping (raw correlation $r \approx 0.0004$)}
\label{fig:house-price-scatter}
\end{figure}

\begin{figure}[H]
\centering
\includegraphics[width=0.85\textwidth]{house_price_bed_bath_box.png}
\caption{Price by bed count (left) and bath count (right), zoomed to the 95th percentile --
both reveal the same outlier-driven distortion in the raw category range before capping}
\label{fig:house-price-bedbath}
\end{figure}

\begin{figure}[H]
\centering
\includegraphics[width=0.55\textwidth]{house_price_state_median.png}
\caption{Median price by state -- confirms \texttt{state} carries real, keep-worthy price
signal (California/Hawaii/D.C./Utah highest, Michigan/Arkansas/West Virginia/Alaska
lowest)}
\label{fig:house-price-state}
\end{figure}

\begin{figure}[H]
\centering
\includegraphics[width=0.62\textwidth]{house_price_corr_heatmap.png}
\caption{Correlation heatmap: numeric listing attributes + price. \texttt{bath} shows the
strongest raw correlation with price ($r=0.39$), followed by \texttt{bed} ($r=0.22$) and
\texttt{house\_size} ($r=0.12$); \texttt{acre\_lot} is weakest ($r=0.02$).}
\label{fig:house-price-corr}
\end{figure}

\subsection{Representation}
\begin{itemize}
    \item \textbf{Scaled numeric} (\texttt{StandardScaler}, fit on train only): capped
    \texttt{bed}, \texttt{bath}, \texttt{acre\_lot}, \texttt{house\_size}, plus
    \texttt{days\_since\_prev\_sale} -- 5 columns.
    \item \textbf{One-hot categorical} (fit on train only): \texttt{status} and
    \texttt{state} -- 59 columns.
    \item \textbf{Frequency-encoded}: \texttt{city} (20,098 distinct values -- too
    high-cardinality for one-hot, but the state-level price spread showed genuine
    sub-state location signal was worth keeping) -- 1 column. \texttt{zip\_code} was
    dropped for the same cardinality reason without the same expected payoff.
\end{itemize}
Resulting shape: $X \in \mathbb{R}^{599,522 \times 65}$.

\subsection{Classical Machine-Learning Models}
\begin{table}[H]
\centering
\begin{tabularx}{\textwidth}{@{}l X@{}}
\toprule
\textbf{Model} & \textbf{Why} \\
\midrule
Linear Regression & Simplest baseline every other model is implicitly measured against. \\
Random Forest & Model deployed in Assignment 02 for this app -- kept for continuity
(\texttt{max\_depth=20}, found necessary in practice: unbounded depth on 480K rows
produced multi-hundred-thousand-node trees and an 8GB+ model file). \\
Gradient Boosting & Sequential residual-correction, structurally different from Random
Forest. \\
\bottomrule
\end{tabularx}
\caption{House Price: 3 classical models and selection rationale}
\end{table}

\subsection{From-Scratch Deep Learning Model}
\textbf{Architecture:} $65 \to 64 \to 32 \to 1$, ReLU hidden layers, \textbf{linear} output
(no activation -- regression, not classification), MSE loss. Wider hidden layers than
diabetes's despite the smaller input, since log-price's relationship to bed/bath/size/
location is plausibly more nonlinear, giving the network more room to compose interactions
before collapsing to one output.

\textbf{Forward propagation:} identical two-layer ReLU structure to diabetes's network, but
the output layer is deliberately \textbf{linear} ($\hat{y}=Z_3$, no sigmoid) -- the one
required departure from the classification networks' math, since squashing to $(0,1)$ would
cap predicted log-price at 1, which is wrong for a target that must range freely.

\textbf{Backpropagation:} the one place the math must differ, not just be copied: $dZ_3 =
2(\hat{y}-y)/n$ is MSE's gradient with respect to a linear output taken directly, with no
sigmoid derivative to cancel against (unlike the classification networks' simpler-looking
$\hat{y}-y$ form). Every layer after that follows the identical
route-back$\to$ReLU-gate$\to$gradient pattern used in the diabetes network. Separately
gradient-checked against a numerical finite-difference estimate (maximum relative error
$\approx 4\times 10^{-9}$) before training on real data.

\textbf{Training:} 200 epochs, full-batch gradient descent, learning rate 0.05 (lower than
diabetes's 0.1, chosen for the wider 64/32 hidden layers' larger accumulated gradient
magnitudes). Loss fell from 152.14 to 1.04, settling by roughly epoch 60
(Figure~\ref{fig:house-price-loss}).

\begin{figure}[H]
\centering
\includegraphics[width=0.65\textwidth]{house_price_loss_curve.png}
\caption{From-scratch DL model: training loss (MSE, log-price space) vs.\ epoch}
\label{fig:house-price-loss}
\end{figure}

\subsection{Evaluation: All 4 Models}
\begin{table}[H]
\centering
\begin{tabular}{lccc}
\toprule
\textbf{Model} & \textbf{RMSE (USD)} & \textbf{MAE (USD)} & \textbf{$R^2$} \\
\midrule
Linear Regression & 894,828 & 271,849 & 0.2599 \\
\textbf{Random Forest} & \textbf{794,082} & \textbf{207,440} & \textbf{0.4172} \\
Gradient Boosting & 900,739 & 244,026 & 0.2501 \\
DL (from scratch) & 1,010,214 & 328,971 & 0.0567 \\
\bottomrule
\end{tabular}
\caption{House Price: 4-model test-set comparison (dollar-space, after $\text{expm1}$
inversion)}
\label{tab:house-price-results}
\end{table}

\begin{figure}[H]
\centering
\includegraphics[width=0.72\textwidth]{house_price_r2_bar.png}
\caption{$R^2$ across all 4 models}
\label{fig:house-price-r2}
\end{figure}

\begin{figure}[H]
\centering
\includegraphics[width=0.8\textwidth]{house_price_rmse_mae_bar.png}
\caption{RMSE / MAE in USD across all 4 models}
\label{fig:house-price-rmse-mae}
\end{figure}

\begin{figure}[H]
\centering
\includegraphics[width=0.62\textwidth]{house_price_pred_vs_actual.png}
\caption{Random Forest: predicted vs.\ actual price, log-log scale (5K-point test-set
sample). Points track the perfect-prediction diagonal reasonably well through the middle
of the price range, with spread opening up at both extremes -- exactly where a
heavy-tailed target is hardest to predict.}
\label{fig:house-price-pred-actual}
\end{figure}

\begin{figure}[H]
\centering
\includegraphics[width=0.75\textwidth]{house_price_residuals.png}
\caption{Random Forest: residual plot in log-price space (5K-point test-set sample).
Residuals center near zero across most of the predicted range with a mild fan-out at the
extremes -- a moderately healthy pattern, not a strong funnel or curve.}
\label{fig:house-price-residuals}
\end{figure}

\begin{figure}[H]
\centering
\includegraphics[width=0.78\textwidth]{house_price_feature_importance.png}
\caption{Top 15 feature importances, Random Forest. \texttt{bath} dominates (0.326),
followed by \texttt{house\_size} (0.212) and \texttt{acre\_lot} (0.138) -- notably,
\texttt{acre\_lot} ranks far higher here than its weak raw correlation with price
suggested, indicating the ensemble captures a nonlinear or interaction effect a simple
Pearson correlation cannot.}
\label{fig:house-price-fi}
\end{figure}

\subsection{Written Comparison}
Random Forest wins clearly on every metric -- $R^2$ 0.417, RMSE \$794,082, MAE \$207,440 --
well ahead of Linear Regression ($R^2$ 0.260), Gradient Boosting ($R^2$ 0.250, sklearn
defaults, untuned), and the from-scratch DL model, which places last ($R^2$ 0.057, RMSE
\$1,010,214, MAE \$328,971).

Figure~\ref{fig:house-price-pred-actual} shows Random Forest's predictions tracking the
perfect-prediction diagonal well through the bulk of the price range, with the expected
spread at the extremes; Figure~\ref{fig:house-price-residuals} shows no strong funnel or
curvature in the residuals, consistent with a moderately-fit, not badly misspecified,
model on a genuinely hard target. Figure~\ref{fig:house-price-fi} adds a further nuance:
Random Forest relies most on \texttt{bath}, but weights \texttt{acre\_lot} far more
heavily than its raw correlation with price would suggest ($r=0.02$ vs.\ an importance of
0.138) -- direct evidence that the ensemble is capturing a nonlinear or
location-interaction effect (lot size mattering differently by state/city) that the
correlation heatmap in Figure~\ref{fig:house-price-corr} alone could not reveal.

In the lecture's framing: the from-scratch network has comparable raw capacity to the tree
ensembles, and its training loss fell smoothly with no divergence, so this is not primarily
an optimization failure. The real driver is the target's shape colliding with this model
family's structure: \texttt{price} is extremely heavy-tailed even after the
$\log(1+y)$ transform (log-price ranges from 0.69 to 19.5, and the EDA found near-billion-
dollar and near-zero listings that are themselves data artifacts), and a from-scratch dense
network with a single unbounded linear output layer has no mechanism analogous to a tree
ensemble's leaf-averaging to naturally cap its response to outlier-heavy regions of the
input space. Random Forest predicts by averaging many trees' leaf values, each fit to a
bootstrap resample -- an implicit regularization against exactly this kind of heavy-tailed,
noisy target that a plain feedforward network trained by unweighted MSE does not get for
free. This mirrors a known, general limitation of small from-scratch networks on tabular
data with heavy-tailed targets, not a bug in the backpropagation math (separately
gradient-checked, Section 4.5) -- a validation-based learning-rate schedule, gradient
clipping, or Huber loss in place of MSE are natural next steps, left out here to keep the
training loop matching the lecture's exact MSE formulation.

% ---------------------------------------------------------
% SECTION 5: APPLICATION 3 - CUSTOMER BEHAVIOUR
% ---------------------------------------------------------
\section{Application 3 --- Customer Behaviour (Flipkart Reviews)}

\subsection{Problem Description}
Binary classification: predict whether a Flipkart product review is a "recommend"
(\texttt{Rate} $\geq 4$) from the review text and product price, on a different e-commerce
platform from Assignment 02's Sephora skincare reviews.
\[
f: \mathbb{R}^{8696} \longrightarrow \{0, 1\} \quad \text{(classical models)}
\]

\subsection{Dataset}
\begin{itemize}
    \item \textbf{Source:} Kaggle dataset \texttt{niraliivaghani/flipkart-dataset},
    \texttt{Dataset.csv} (saved locally as \texttt{flipkart\_reviews.csv}).
    \item \textbf{Rows:} 363,261 raw $\to$ 363,239 after dropping 22 structurally corrupted
    rows (0.006\%, a CSV parsing artifact where the header row and a handful of blank rows
    appear as data) $\to$ 290,591 train / 72,648 test (80/20 stratified split).
    \item \textbf{Growth vs.\ Assignment 02:} 3.1$\times$ the prior 116,264-row Sephora
    dataset, and a deliberately different platform (general merchandise, not skincare) to
    remove any ambiguity about whether the growth requirement was satisfied by genuinely
    new data.
    \item \textbf{Target:} \texttt{is\_recommended} $= \mathbb{1}[\texttt{Rate} \geq 4]$.
\end{itemize}

\subsection{Data Understanding and Cleaning}
The raw \texttt{Rate} distribution is heavily skewed toward high scores: 5-star is the
largest bucket (203,832 reviews, 56.1\%), followed by 4-star (74,252, 20.4\%); 1-star
(40,077), 3-star (32,098), and 2-star (12,980) are all smaller
(Figure~\ref{fig:cb-rate}). Collapsing to the binary target gives 278,084 positive vs.\
85,155 negative reviews -- a 76.5\%/23.5\% split. Median review length is longest for the
most negative reviews (15 characters at 1- and 2-star) and shortest for neutral 3-star
reviews (9 characters), rising again for positive reviews (11--13 characters) -- a weak,
non-monotonic signal, which is why \texttt{Review} is used as full TF-IDF content rather
than reduced to a length feature (Figure~\ref{fig:cb-review-length}). The raw \texttt{Price}
column decodes into strings of only 3 distinct lengths because the rupee (INR) symbol was
corrupted during download into stray bytes; \texttt{clean\_price()} strips everything
except digits/commas/decimal points via regex before parsing, recovering a standard
right-skewed price distribution (median INR 1,699, mean INR 4,582, max INR 86,990;
Figure~\ref{fig:cb-price-cleaning}). \texttt{Product\_name} has only 1,175 distinct values
across 363K+ reviews, dominated by a handful of high-traffic listings
(Figures~\ref{fig:cb-top-products-1}--\ref{fig:cb-top-products-2}) -- confirming the
decision to drop it as too high-cardinality and skewed to use as a clean category.

\begin{figure}[H]
\centering
\includegraphics[width=0.85\textwidth]{cb_rate_balance.png}
\caption{Raw \texttt{Rate} distribution 1--5 (left) and derived \texttt{is\_recommended}
class balance (right)}
\label{fig:cb-rate}
\end{figure}

\begin{figure}[H]
\centering
\includegraphics[width=0.6\textwidth]{cb_review_length.png}
\caption{Review character length by star rating}
\label{fig:cb-review-length}
\end{figure}

\begin{figure}[H]
\centering
\includegraphics[width=0.85\textwidth]{cb_price_cleaning.png}
\caption{Raw \texttt{Price} string length before cleaning (left, 3 discrete lengths from
the rupee-symbol corruption artifact) vs.\ parsed numeric price after cleaning, clipped at the 99th
percentile for display (right)}
\label{fig:cb-price-cleaning}
\end{figure}

\begin{figure}[H]
\centering
\includegraphics[width=0.85\textwidth]{cb_top_products_part1.png}
\caption{Top 20 products by review count -- rank 1--10}
\label{fig:cb-top-products-1}
\end{figure}

\begin{figure}[H]
\centering
\includegraphics[width=0.85\textwidth]{cb_top_products_part2.png}
\caption{Top 20 products by review count -- rank 11--20}
\label{fig:cb-top-products-2}
\end{figure}

The raw word vocabulary itself already separates the two classes before any modeling:
the 15 most frequent words in negative (\texttt{Rate} $< 4$) reviews include
\textit{waste}, \textit{fair}, and \textit{money} alongside neutral terms like
\textit{product}, while positive (\texttt{Rate} $\geq 4$) reviews are dominated by
\textit{awesome}, \textit{terrific}, and \textit{good} (Figure~\ref{fig:cb-word-freq}) --
a useful sanity check, run before any TF-IDF weighting, that the text genuinely carries
class-discriminating signal.

\begin{figure}[H]
\centering
\includegraphics[width=0.95\textwidth]{cb_word_frequency.png}
\caption{Top 15 words by raw frequency, negative reviews (left) vs.\ positive reviews
(right), stop-words removed}
\label{fig:cb-word-freq}
\end{figure}

\subsection{Representation}
\begin{itemize}
    \item \textbf{Review text} $\to$ TF-IDF, 1--2 gram, two separate vocabularies: the full
    8,695-term vocabulary for the classical models (sklearn's sparse linear models handle
    this natively even at 364K rows), capped at \texttt{max\_features=2000} for the DL
    model. The cap was reached empirically, not chosen from theory alone: an original
    5,000-feature plan caused real memory exhaustion (\(\sim\)11.7GB for the dense training
    array alone at float64) when the sparse TF-IDF matrix had to be densified for plain
    NumPy matrix multiplication; 2,000 features at \texttt{float32} keeps the dense array
    to \(\sim\)2.3GB.
    \item \textbf{Price} (cleaned): \texttt{StandardScaler}, fit on train only.
    \item \textbf{Product\_name}: dropped (see Data Understanding above).
\end{itemize}
Resulting shapes: classical models $X \in \mathbb{R}^{363,239 \times 8,696}$ (sparse); DL
model $X \in \mathbb{R}^{363,239 \times 2,001}$ (densified to \texttt{float32}).

\subsection{Classical Machine-Learning Models}
\begin{table}[H]
\centering
\begin{tabularx}{\textwidth}{@{}l X@{}}
\toprule
\textbf{Model} & \textbf{Why} \\
\midrule
Logistic Regression & Model deployed in Assignment 02 for this app -- kept for continuity. \\
Linear SVM (\texttt{SGDClassifier}, hinge loss) & Max-margin objective well-suited to
high-dimensional sparse text; \texttt{LinearSVC} itself failed to converge in 30+ minutes
on 291K rows, so SGD's incremental solver was used instead to reach the same objective. \\
Multinomial Naive Bayes & Cheap, fast text baseline; trained on text alone (Naive Bayes
requires non-negative count-like input, which scaled \texttt{Price} would violate). \\
\bottomrule
\end{tabularx}
\caption{Customer Behaviour: 3 classical models and selection rationale}
\end{table}

\subsection{From-Scratch Deep Learning Model}
\textbf{Architecture:} $2{,}001 \to 128 \to 32 \to 1$, ReLU hidden layers, sigmoid output,
binary cross-entropy loss. The widest first hidden layer of the three applications, since
the sparse, high-dimensional TF-IDF input needs more first-layer capacity to compress
usefully -- the same narrowing pattern the lecture material illustrates for
high-dimensional inputs.

\textbf{Forward propagation:} layer 1 compresses the \(\sim\)2,001-wide, mostly-zero input
(2,000 TF-IDF term weights + 1 scaled price) down to 128 dense features, each a learned
combination of TF-IDF terms; layer 2 recombines those into 32 higher-level features; the
output layer applies sigmoid, exactly as in the diabetes network, to produce a
recommend-probability.

\textbf{Backpropagation:} identical structure to the diabetes network ($dZ_3 = \hat{y}-y$,
same sigmoid+BCE cancellation), walked through once more here because the input is TF-IDF
rather than tabular survey fields -- \texttt{dW1 = X}$^{\top}$\texttt{dZ1} ends up a
$(2001, 128)$ matrix, one gradient for every TF-IDF term weight feeding into every
first-layer neuron, computed in a single matrix multiplication. Separately gradient-checked
(maximum relative error $\approx 5\times 10^{-8}$) before training on real data.

\textbf{Training:} 200 epochs, full-batch gradient descent, learning rate 0.1, inputs cast
to \texttt{float32} to keep the densified training matrix within memory. Loss fell from
0.698 to 0.382 (Figure~\ref{fig:cb-loss}) -- decreasing, but far less steeply than either
other application's DL model.

\begin{figure}[H]
\centering
\includegraphics[width=0.65\textwidth]{cb_loss_curve.png}
\caption{From-scratch DL model: training loss (binary cross-entropy) vs.\ epoch}
\label{fig:cb-loss}
\end{figure}

\subsection{Evaluation: All 4 Models}
\begin{table}[H]
\centering
\begin{tabular}{lcccc}
\toprule
\textbf{Model} & \textbf{Accuracy} & \textbf{Precision} & \textbf{Recall} & \textbf{F1} \\
\midrule
\textbf{Logistic Regression} & \textbf{0.9943} & \textbf{0.9965} & 0.9960 & \textbf{0.9963} \\
Linear SVM & 0.9940 & 0.9952 & \textbf{0.9969} & 0.9961 \\
Multinomial NB & 0.9674 & 0.9626 & 0.9961 & 0.9790 \\
DL (from scratch) & 0.7656 & 0.7656 & 1.0000 & 0.8672 \\
\bottomrule
\end{tabular}
\caption{Customer Behaviour: 4-model test-set comparison}
\label{tab:cb-results}
\end{table}

\begin{figure}[H]
\centering
\includegraphics[width=0.72\textwidth]{cb_accuracy_bar.png}
\caption{Accuracy across all 4 models}
\label{fig:cb-acc}
\end{figure}

\begin{figure}[H]
\centering
\includegraphics[width=0.8\textwidth]{cb_prf1_bar.png}
\caption{Precision / Recall / F1 across all 4 models}
\label{fig:cb-prf1}
\end{figure}

\begin{figure}[H]
\centering
\includegraphics[width=0.55\textwidth]{cb_confusion_matrix.png}
\caption{From-scratch DL model: confusion matrix on the test set -- every row was
predicted class 1, including all 17,031 true negatives}
\label{fig:cb-cm}
\end{figure}

\begin{figure}[H]
\centering
\includegraphics[width=0.72\textwidth]{cb_roc_curves.png}
\caption{ROC curves across all 4 models. The from-scratch DL model's AUC (0.9974) is
essentially tied with all three classical models (0.9974--0.9981) -- a $\sim$0.0007 gap,
nowhere close to the $\sim$23-point Accuracy gap in Table~\ref{tab:cb-results}.}
\label{fig:cb-roc}
\end{figure}

\begin{figure}[H]
\centering
\includegraphics[width=0.78\textwidth]{cb_top_terms.png}
\caption{Strongest TF-IDF terms for/against ``recommend'' by Logistic Regression
coefficient. Positive terms (\textit{fabulous}, \textit{brilliant}, \textit{wonderful},
\textit{awesome}) and negative terms (\textit{horrible}, \textit{worthless},
\textit{terrible}, \textit{poor}) are exactly the sentiment vocabulary a human reader
would expect.}
\label{fig:cb-top-terms}
\end{figure}

\subsection{Written Comparison}
The 3 classical models all score above 0.99 Accuracy -- Logistic Regression highest
(0.9943, F1 0.9963), Linear SVM close behind (0.9940, F1 0.9961), Multinomial Naive Bayes
lowest of the three but still strong (0.9674, F1 0.9790). The from-scratch DL model trails
far behind at 0.7656 Accuracy, and its confusion matrix ($[[0,17031],[0,55617]]$) shows
exactly why: Recall is a perfect 1.0 but Precision equals Accuracy exactly, because the
model predicted class 1 ("recommend") for every single test row, including all 17,031 true
negatives. It degenerated to the majority-class baseline rather than learning a real
decision boundary.

The ROC curves in Figure~\ref{fig:cb-roc} tell a strikingly different story from the
confusion matrix alone: the DL model's AUC (0.9974) is essentially tied with all three
classical models (0.9974--0.9981). This is the clearest evidence in this application that
the DL model's collapse to always-predict-positive is purely a threshold artifact, not a
representation or learning failure -- the network assigns almost exactly the same
relative ranking to reviews as the classical models do, it simply never crosses 0.5 for
the minority class under unweighted BCE. Figure~\ref{fig:cb-top-terms} reinforces this
from a different angle: Logistic Regression's strongest terms are exactly the sentiment
words a human would expect, confirming the TF-IDF representation itself carries strong
signal -- signal the DL model's own architecture (with its capped 2,000-term vocabulary)
had less of to work with in the first place.

In the lecture's framing, this is neither an overfitting problem nor primarily a
representation-capacity problem in the usual sense -- training loss did decrease over 200
epochs, just not far enough to move predictions off the majority-class default, and the
ROC curve confirms the ranking signal was there all along. Two compounding causes explain
the threshold failure specifically. First, the DL model's input vocabulary is
deliberately capped at 2,000 TF-IDF terms against the classical models' full 8,695-term
vocabulary -- a more than 4$\times$ reduction in the exact signal (specific words
correlating with negative reviews) that made the classical linear models so effective;
TF-IDF-as-input to a dense network throws away the sparse, high-dimensional structure
that Logistic Regression and Linear SVM exploit natively. Second, the 76.5\%/23.5\% class
imbalance in \texttt{is\_recommended} gives unweighted binary cross-entropy a cheap,
loss-reducing shortcut -- predict positive for everyone and start already at 76.5\%
"accuracy" from the loss function's perspective -- and 200 epochs of plain full-batch
gradient descent did not escape that basin. This is the same class-imbalance dynamic that
hurt the diabetes DL model's Recall (Section 3.7), compounded here by the capped
representation. Both limitations are genuine and explainable rather than implementation
bugs (separately gradient-checked, Section 5.5), and both point to the same fix used for
diabetes: class-weighted loss and/or a tuned decision threshold, left out here to keep the
training loop matching the lecture's exact unweighted BCE formula.

% ---------------------------------------------------------
% SECTION 6: CROSS-APPLICATION COMPARISON
% ---------------------------------------------------------
\section{Cross-Application Comparison}

\begin{table}[H]
\centering
\begin{tabular}{lcccc}
\toprule
\textbf{Application} & \textbf{Classical 1} & \textbf{Classical 2} & \textbf{Classical 3} &
\textbf{DL (from scratch)} \\
\midrule
Diabetes (Accuracy) & LogReg: 0.850 & RF: 0.849 & \textbf{GB: 0.852} & 0.845 \\
House Price ($R^2$) & LinReg: 0.260 & \textbf{RF: 0.417} & GB: 0.250 & 0.057 \\
Customer Behaviour (Accuracy) & \textbf{LogReg: 0.994} & SVM: 0.994 & NB: 0.967 & 0.766 \\
\bottomrule
\end{tabular}
\caption{Primary metric across all 3 applications $\times$ 4 models (bold = winner per
application)}
\label{tab:cross-app}
\end{table}

\begin{table}[H]
\centering
\small
\renewcommand{\arraystretch}{1.35}
\begin{tabularx}{\textwidth}{@{}l X X X@{}}
\toprule
\textbf{Dimension} & \textbf{Diabetes} & \textbf{House Price} & \textbf{Customer
Behaviour} \\
\midrule
Task & Binary classification & Regression & Binary classification \\
Rows (raw) & 815,711 & 2,226,382 (600K sampled) & 363,261 \\
Growth vs.\ A02 & 3.2$\times$ & 2.5$\times$ & 3.1$\times$ \\
Representation & Dense matrix, $\mathbb{R}^{52}$ & Dense matrix, $\mathbb{R}^{65}$ &
Sparse TF-IDF + tabular, $\mathbb{R}^{8,696}$ \\
DL architecture & $52\to32\to16\to1$ & $65\to64\to32\to1$ (linear out) &
$2001\to128\to32\to1$ \\
DL primary-metric rank & 4th of 4 & 4th of 4 & 4th of 4 \\
DL failure mode & Class-imbalance-driven low Recall & Heavy-tailed target, unbounded
output & Class-imbalance-driven majority-class collapse \\
\bottomrule
\end{tabularx}
\caption{Systemic comparison across all three applications}
\label{tab:cross-app-2}
\end{table}

% ---------------------------------------------------------
% SECTION 7: DISCUSSION
% ---------------------------------------------------------
\section{Discussion: What Changed From Assignment 02}

\textbf{Diabetes.} The dataset grew 3.2$\times$ by adding a second BRFSS survey year
(2019+2020) rather than a new source, keeping the same feature family so the comparison
against Assignment 02's classical results stays meaningful. The genuinely new ingredient
this round is the from-scratch DL model, which reached a competitive Accuracy (0.845, only
0.7 points behind the best classical model) but a clearly worse Recall (0.028 vs.\ 0.191)
-- a direct consequence of training on unweighted BCE over an imbalanced target with plain
gradient descent, discussed in Section 3.7.

\textbf{House Price.} The dataset grew 2.5$\times$ and switched markets entirely (US
listings instead of Vietnamese listings), which changed both scale (USD, not million VND)
and the specific outlier patterns encountered (32-bit-overflow-style corruption in price,
house size, bed, and bath counts, cleaned by percentile capping rather than the
missing-indicator approach the house-price application's own $\sim$70--84\% structural
missingness needed in Assignment 02). Training a pure-NumPy loop on 480K rows and 65
features for 200 epochs was noticeably heavier than the other two applications' networks, though still
tractable on a single machine without a GPU. The from-scratch DL model did not beat the
classical models here -- a legitimate finding, not a shortcoming to gloss over: a small
feedforward network's unbounded linear output has no protection against a heavy-tailed
target the way a tree ensemble's leaf-averaging does (Section 4.7).

\textbf{Customer Behaviour.} The dataset grew 3.1$\times$ and deliberately switched
e-commerce platforms (Flipkart general merchandise instead of Sephora skincare) to remove
any ambiguity about whether reusing the same source would count as genuine growth. What got
harder here specifically was memory, not just training time: the DL model's dense TF-IDF
input had to be capped from a planned 5,000 features down to 2,000 (and cast to
\texttt{float32}) after an earlier attempt exhausted available RAM and swapped -- a direct,
concrete cost of training a from-scratch network on text without a framework's sparse
autodiff support. The from-scratch DL model did not win here either, and by the widest
margin of the three applications: it collapsed to a majority-class predictor, a known,
explainable limitation of unweighted BCE plus a severely capped vocabulary rather than an
implementation defect (Section 5.7).

\textbf{Synthesis.} Across all three applications the from-scratch deep learning model
never won its task's primary metric. This is a genuine, useful finding rather than a
disappointing one: it demonstrates two well-known, general limitations of small
from-scratch networks trained with plain, unweighted, full-batch gradient descent --
vulnerability to class imbalance (diabetes, customer behaviour) and lack of the implicit
regularization tree ensembles get from bootstrap averaging on heavy-tailed tabular targets
(house price) -- limitations that are well-documented in the deep learning literature and
not specific to this implementation. Every backpropagation derivation was independently
verified by numerical gradient-checking before being trained on real data, confirming the
math itself is correct; what differs from the classical models is the training procedure's
robustness to imbalance and outliers, not the network's mechanical correctness.

% ---------------------------------------------------------
% SECTION 8: CONCLUSION
% ---------------------------------------------------------
\section{Conclusion}

This assignment extended three intelligent systems with a deep learning model built
entirely from scratch in NumPy -- forward propagation, loss, backpropagation via the chain
rule, and gradient descent all derived and coded by hand, with every derivation separately
verified against numerical gradient checks. The from-scratch models did not outperform the
classical scikit-learn baselines on any of the three applications, but the reasons why are
themselves the main lesson of this assignment: representation and training-procedure
choices (vocabulary size, class weighting, output-layer activation matched to the target's
support) matter as much as, or more than, a network's raw capacity, and a correctly
implemented network can still lose to a well-suited classical algorithm when its training
procedure is not adapted to the data's imbalance or heavy tails. Understanding \emph{why}
each from-scratch model fell short -- rather than only reporting that it did -- is what
this exercise in building deep learning from first principles was for.

% ---------------------------------------------------------
% SECTION 9: REPRODUCIBILITY
% ---------------------------------------------------------
\section{Reproducibility}

\begin{itemize}
    \item \textbf{Environment:} Python 3.x, macOS.
    \item \textbf{Core libraries:} \texttt{numpy}, \texttt{pandas}, \texttt{scikit-learn},
    \texttt{matplotlib}, \texttt{seaborn}, \texttt{joblib}. No \texttt{tensorflow},
    \texttt{torch}, or \texttt{keras} anywhere in the from-scratch DL implementation.
    \item \textbf{Determinism:} \texttt{RANDOM\_SEED = 42} throughout -- the 600K-row
    house-price sample, every train/test split, every sklearn estimator, and every
    from-scratch network's initial weights are all seeded from this same value.
    \item \textbf{No-leakage discipline:} every scaler, encoder, and imputation statistic
    is fit on the training split only and reused (never refit) on the test split, matching
    Assignment 02's rule.
    \item \textbf{Re-running a notebook end-to-end:}
\begin{lstlisting}
jupyter nbconvert --to notebook --execute --inplace \
    diabetes/notebook/diabetes.ipynb
jupyter nbconvert --to notebook --execute --inplace \
    house_price/notebook/house_price.ipynb
jupyter nbconvert --to notebook --execute --inplace \
    customer_behaviour/notebook/customer_behaviour.ipynb
\end{lstlisting}
    \item \textbf{Artifacts produced per app} (in \texttt{<app>/model/}):
    \texttt{model\_pipeline.joblib} (best classical model by primary metric),
    \texttt{dl\_weights.npz} (from-scratch DL model's 6 weight/bias arrays),
    \texttt{feature\_names.joblib}, \texttt{input\_schema.json}. No API, web, or mobile
    deployment layer exists for this assignment -- both source lecture decks stop at the
    modeling pipeline and never mention deployment.
\end{itemize}

\end{document}
